\chapter{Many Body Theory}

Many body theory is a powerful framework for understanding the behavior of systems with a large number of interacting particles. In condensed matter physics, it is essential for describing phenomena such as superconductivity, magnetism, and the electronic properties of materials. This theory was developed to address the complexities that arise when dealing with systems of many particles, where the perturbation theory of non-interacting particles fails to capture the essential physics. This theory was developed in the 1957-1960s by pioneers such as John Bardeen, Leon Cooper, and Robert Schrieffer, who formulated the BCS theory of superconductivity. The many body theory has since been extended and refined to include various interactions and correlations, making it a cornerstone of modern condensed matter physics.

In this chapter, we will explore the fundamental concepts and techniques of many body theory, including Green's functions, Feynman diagrams, and the random phase approximation. We will also discuss how these tools are used to analyze the properties of interacting electron systems and to understand phenomena such as collective excitations and phase transitions.

\subsubsection{Time Dependencies in Quantum Theories}

Firstly we need to speak about time dependencies in quantum theories. Since we are dealing with particles being created and destroyed after some kind of propagation, we need to understand how to deal with time dependencies in quantum mechanics.

\paragraph{Schrodinger picture.} In the Schrodinger picture, the state vectors evolve in time while the operators are time-independent. The time evolution of a state vector \(\ket{\psi}\) is governed by the Schrodinger equation:
\[
    \dots
\]

\paragraph{Heisenberg picture.} In the Heisenberg picture, the operators evolve in time while the state vectors are time-independent. The time evolution of an operator \(A\) is given by:
\[
    \dots
\]

we can thus derive an equation of motion for the operators in the Heisenberg picture, which is known as the Heisenberg equation of motion:
\[
    \hat{A}_H(t) = e^{i \hat{H}t} A_S e^{-i \hat{H}t},
\]
\[
    \begin{aligned}
        \frac{\mathrm{d}}{\mathrm{d} t} \hat{A}_H(t) & = \frac{\mathrm{d}}{\mathrm{d} t} \left( e^{i \hat{H}t} A_S e^{-i \hat{H}t} \right)  = i \hat{H} e^{i \hat{H}t} A_S e^{-i \hat{H}t} - e^{i \hat{H}t} A_S i \hat{H} e^{-i \hat{H}t} \\
                                                     & = i \hat{H} \hat{A}_H(t) - i \hat{A}_H(t) \hat{H} = i [\hat{H}, \hat{A}_H(t)].
    \end{aligned}
\]

\paragraph{Interaction (Dirac) picture.}
Sometimes we have an hamiltonian which is composed of a time independent part and a time independent part, such as:
\[
    \hat{H} = \hat{H}_0 + \hat{V}(t).
\]
In this case, we can define the interaction picture, where both the state vectors and the operators have time dependence. The operators in the interaction picture evolve according to the free Hamiltonian \(\hat{H}_0\) (or the quadratic part of the Hamiltonian), while the state vectors evolve according to the interaction part \(\hat{V}(t)\):
\[
    A_I(t) = e^{i \hat{H}_0 t} A_S e^{-i \hat{H}_0 t},
\]
\[
    \ket{\psi_I(t)} = e^{i \hat{H}_0 t} \ket{\psi_S} e^{-i \hat{H}_0 t}.
\]
The interaction picture is particularly useful in many body theory, as it allows us to treat the interactions perturbatively while still accounting for the time evolution of the system. It makes evident the connection between the free and interacting theories, and is often used in the formulation of Feynman diagrams and perturbation theory.

\paragraph{Heisenberg - Dirac bridge.}
We can also derive a connection between the Heisenberg and Dirac pictures. Starting from the definition of the Heisenberg picture, we can add identity operators to express the Heisenberg operator in terms of the interaction picture:
\[
    \begin{aligned}
        A_H(t) = e^{i \hat{H} t} A_S e^{-i \hat{H} t} & = e^{i \hat{H} t} e^{i \hat{H}_0(t) t} e^{-i \hat{H}_0(t) t} A_S e^{-i \hat{H}_0(t) t} e^{i \hat{H}_0(t) t} e^{-i \hat{H}t} \\
                                                      & = U^\dagger(t,0) A_I(t) U(t,0),
    \end{aligned}
\]
where we have defined the time evolution operator in the interaction picture as:
\[
    U(t,0) = e^{i \hat{H}_0 t} e^{-i \hat{H} t} = e^{i \hat{H}_0 t} e^{-i (\hat{H}_0 + \hat{V}(t)) t} = e^{-i \int_0^t \mathrm{d} t' \hat{V}_I(t')}.
\]
This evolution operator contains only the interaction part of the Hamiltonian, and this highlights the fact that the interaction picture is designed to evolve the state vectors according to the interaction while keeping the operators evolving according to the free Hamiltonian.

Pay attention:
\[
    e^{i \hat{H}_0 t} e^{-i \hat{H} t} \neq e^{i \hat{V} t},
\]
and this is the kernel of many body theory, since the interaction part of the Hamiltonian does not commute with the free part normally, and thus we cannot simply factorize the time evolution operator into a product of exponentials. This non-commutativity is what gives rise to the rich physics of interacting systems, and it is what makes many body theory both challenging and fascinating.

Let us take the partial time derivative of the evolution operator:
\[
    \begin{aligned}
        \frac{\partial}{\partial t} \hat{U}(t,0) & = \dots                         \\
                                                 & = \dots                         \\
                                                 & = -i \hat{V}_I(t) \hat{U}(t,0).
    \end{aligned}
\]
Now the question is: how do we solve this equation? We can formally integrate it to get:
\[
    \hat{U}(t,0) - \hat{U}(0,0) = - i \int_0^t \mathrm{d} t' \hat{V}_I(t') \hat{U}(t',0).
\]
This is an integral equation for the evolution operator, and it can be solved iteratively by substituting the expression for \(\hat{U}(t',0)\) back into the integral. This leads to a series expansion known as the \textbf{Dyson series}:
\[
    \hat{U}(t,0) = 1 - i \int_0^t \mathrm{d} t' \hat{V}_I(t') + (-i)^2 \int_0^t \mathrm{d} t' \int_0^{t'} \mathrm{d} t'' \hat{V}_I(t') \hat{V}_I(t'') + \dots
\]
This series expansion allows us to compute the evolution operator perturbatively, and it forms the basis for many of the techniques used in many body theory, such as Feynman diagrams and perturbation theory. Each term in the Dyson series corresponds to a different order of interaction, and by summing over all possible interactions, we can capture the full dynamics of the system.

If we take a look at the domain of integration in the Dyson series, we can see that it is ordered in time, with \(t' > t''\). This time ordering is crucial for ensuring causality in the theory, and it is often represented using the time-ordering operator \(\hat{T}\), which arranges operators in chronological order. The Dyson series can thus be rewritten using the time-ordering operator knowing that
\[
    \begin{aligned}
        \int_0^t \mathrm{d} t' \hat{V}_I(t') \int_0^{t'} \mathrm{d} t'' \hat{V}_I(t'') & = \frac{1}{2} \int_0^t \mathrm{d} t' V_I(t') \int_0^{t^{\prime}} \mathrm{d} t'' \hat{V}_I(t'') + \frac{1}{2} \int_0^t \mathrm{d} t'' \hat{V}_I(t'') \int_0^{t''} \mathrm{d} t' \hat{V}_I(t') \\
                                                                                       &
    \end{aligned}
\]
where now we can introduce the Heaviside step function \(\theta(t)\) to enforce the time ordering of the operators (changing accordingly the integration domain) as
\[
    =\frac{1}{2} \int_0^t \mathrm{d} t' V_I(t') \int_0^{t} \mathrm{d} t'' \hat{V}_I(t'') \theta(t' - t'') + \frac{1}{2} \int_0^t \mathrm{d} t'' \hat{V}_I(t'') \int_0^{t} \mathrm{d} t' \hat{V}_I(t') \theta(t'' - t').
\]
Thus now we can write the Dyson series in a compact form using the time-ordering operator:
\[
    =\frac{1}{2} \int_0^t \mathrm{d} t' \int_0^t \mathrm{d} t'' \left( \hat{V}_I(t') \hat{V}_I(t'') + \hat{V}_I(t'') \hat{V}_I(t') \right) = \frac{1}{2} \int_0^t \mathrm{d} t' \int_0^t \mathrm{d} t'' \hat{T}_t \left( \hat{V}_I(t') \hat{V}_I(t'') \right),
\]
where we have used the definition of the time-ordering operator to combine the two terms into a single integral: the term with the smaller time argument is placed to the right of the term with the larger time argument. Let us make a quick example.

\begin{example}[Time ordering of n operators]
    Consider the case of three operators \(\hat{V}_I(t_1)\), \(\hat{V}_I(t_2)\), and \(\hat{V}_I(t_3)\). The time-ordered product of these operators is given by all the combinations of the operators ordered according to their time arguments:
    \[
        \hat{T} \left( \hat{V}_I(t_1) \hat{V}_I(t_2) \hat{V}_I(t_3) \right) = \frac{1}{3!} \left( V_I(t_1) \hat{V}_I(t_2) \hat{V}_I(t_3) \theta(t_1 - t_2) \theta(t_2 - t_3) + \dots \right).
    \]
    This means that in the time-ordered product, the operator with the latest time argument (in this case, \(t_3\)) is placed to the left, while the operator with the earliest time argument (in this case, \(t_1\)) is placed to the right. This ordering ensures that the physical processes described by these operators occur in the correct chronological sequence.
\end{example}

With this new formal way of expressing the ordered integrals, we can now write the Dyson series in a compact form as
\[
    U(t,0) = \sum_{n=0}^\infty \frac{(-i)^n}{n!} \int_0^t \mathrm{d} t_1 \dots \int_0^t \mathrm{d} t_n \hat{T}_t \left( \hat{V}_I(t_1) \cdots \hat{V}_I(t_n) \right).
\]
This expression can be recognized as the expansion of the exponential function, and thus we can write the evolution operator formally in a compact form as
\[
    \hat{U}(t,0) = \hat{T} \exp\left( -i \int_0^t \mathrm{d} t' \hat{V}_I(t') \right).
\]
This compact expression encapsulates the entire perturbative expansion of the evolution operator, and it is a fundamental result in many body theory, even if in practice we never compute the evolution operator in this form, but rather we use the Dyson series expansion to compute specific terms in the perturbation theory (i.e. we stop at a certain order in the expansion).

\section{Green's Functions}

In this tractation, we have to distinguish the theory of bosons and fermions, since the statistics of the particles will affect the properties of the Green's functions. For simplicity, we will focus on the case of bosons, but the generalization to fermions is straightforward and will be discussed at the end of this section.

\subsection{Single Particle Green's Function}

The single particle Green's function, also known as the propagator, is a fundamental quantity in many body theory that describes the propagation of a single particle in an interacting system. It is defined as the time-ordered expectation value of the field operators:\footnote{For fermions it will be the same but with anticommutators.}
\[
    \mathcal{G}^R(\mathbf{r}, \sigma, t;\,\mathbf{r}', \sigma', t') = -i \theta(t - t') \langle [\hat{\Psi}_{\sigma} (\mathbf{r}, t),\, \hat{\psi}_{\sigma'}^\dagger(\mathbf{r}', t')] \rangle.
\]
This Green's function encodes information about the energy spectrum, the density of states, and the response of the system to external perturbations. It can be used to compute various physical quantities, such as the spectral function, the self-energy, and the effective mass of quasiparticles.

Let us study the formula introduced above:
\begin{itemize}
    \item The field operators \(\hat{\Psi}_{\sigma} (\mathbf{r}, t)\) and \(\hat{\psi}_{\sigma'}^\dagger(\mathbf{r}', t')\) are the annihilation and creation operators for particles with spin \(\sigma\) at position \(\mathbf{r}\) and time \(t\), and for particles with spin \(\sigma'\) at position \(\mathbf{r}'\) and time \(t'\), respectively. Since they depend on time, they are in the Heisenberg picture, and thus they evolve according to the full Hamiltonian of the system, including interactions.
    \item The ``\(^R\)'' superscript denotes that this is the retarded Green's function, which means that it describes the response of the system to a perturbation applied at time \(t'\) and observed at a later time \(t\). The retarded Green's function is causal, meaning that it vanishes for \(t < t'\), which is enforced by the Heaviside step function \(\theta(t - t')\).
    \item The expectation value \(\langle \cdots \rangle\) is taken with respect to the ground state of the system, or more generally, with respect to the thermal equilibrium state at a given temperature. This is a \textbf{thermal average}:
          \[
              \langle \cdots \rangle = \frac{1}{Z} \mathrm{Tr} \left( e^{-\beta H} \cdots \right) = \frac{1}{Z} \sum_{n} \bra{n} e^{-\beta H} \cdots \ket{n} = \frac{1}{Z} \sum_{n} e^{- \beta E_n} \bra{n} \cdots \ket{n},
          \]
          where \(Z\) is the partition function and \(\beta\) is proportional to the inverse of temperature. We usually don't know the exact spectrum of the complete Hamiltonian, indeed usually only the energy of the ground state, i.e. set to 0, contributes significantly to the sum, and thus we can approximate the thermal average as
          \[
              \langle \cdots \rangle \approx \bra{0} \cdots \ket{0},
          \]
          where \(\ket{0}\) is the ground state of the system. This approximation is valid at low temperatures, where the contributions from excited states are negligible compared to the ground state.
\end{itemize}

Let us compute the Green's function at zero temperature for fermions, where we can use the approximation of the thermal average as the expectation value in the ground state \(\ket{GS}\):
\[
    \mathcal{G}^R(\mathbf{r}, \sigma, t;\,\mathbf{r}', \sigma', t') = -i \theta(t - t') \bra{GS} \{\hat{\Psi}_{\sigma} (\mathbf{r}, t),\, \hat{\Psi}_{\sigma'}^\dagger(\mathbf{r}', t') \} \ket{GS},
\]
where we are creating a particle at time \(t'\) and position \(\mathbf{r}'\) and then we are annihilating it at a later time \(t\) and position \(\mathbf{r}\) (retarded GF). This Green's function describes the amplitude for a particle to propagate from \((\mathbf{r}', t')\) to \((\mathbf{r}, t)\) in the presence of interactions. It contains information about the energy spectrum of the system, the density of states, and the response to external perturbations.

This objects is foundamentally an amplitude for a particle to propagate from one point in space and time to another, without taking into account all the infinite possible interactions that can happen in between, and thus it is a very useful quantity to understand the behavior of the system. It is also important to note that the Green's function measures the probability to destroy the previously created particle, in exactely the same state in which it was created.

We are using the commutator and not only the product of the two operators because we want to ensure that the Green's function is properly symmetrized (or antisymmetrized for fermions) and to ensure causality: it can be showed that the commutator (or anticommutator) is the interesting expectation value that captures the correct physical behavior of the system and of the operators, but we will not spend time on this point, since it is not crucial for the understanding of the many body theory.

Remember the structure of the field operators as a sum over quantum states:
\[
    \hat{\Psi}_{\sigma} (\mathbf{r}, t) = \sum_{\nu} \phi_{\nu}(\mathbf{r}) \hat{c}_{\alpha, \sigma}(t) = \sum_{k} \frac{1}{\sqrt{\Omega}} e^{i \mathbf{k} \cdot \mathbf{r}} \hat{c}_{\mathbf{k}, \sigma}(t),
\]
which can be inserted into the definition of the Green's function to express it in terms of the creation and annihilation operators in momentum space of the first quantization:
\[
    \begin{aligned}
        \mathcal{G}^R(\mathbf{r}, \sigma, t;\,\mathbf{r}', \sigma', t') & = -i \theta(t - t') \bra{GS} \{\hat{\Psi}_{\sigma} (\mathbf{r}, t),\, \hat{\Psi}_{\sigma'}^\dagger(\mathbf{r}', t') \} \ket{GS} =                                                                                      \\
                                                                        & = -i \theta(t - t') \bra{GS} \left\{ \sum_{k} \phi_{\nu}(\mathbf{r}) \hat{c}_{\mathbf{k}, \sigma}(t),\, \sum_{k'} \phi_{\nu'}(\mathbf{r}') \hat{c}_{\mathbf{k}', \sigma'}^\dagger(t')  \right\}  \ket{GS} =            \\
                                                                        & = -i \theta(t - t')  e^{i (\mathbf{k} \cdot \mathbf{r} -\mathbf{k}' \cdot \mathbf{r}')}  e^{-i (\epsilon_k t -\epsilon_{k'} t')}  = -i  e^{-i (\epsilon_k (t-t'))}  e^{i (\mathbf{k} \cdot (\mathbf{r}-\mathbf{r}'))},
    \end{aligned}
\]\TODO{Write correct ecxpression.}
which in the momentum space can be expressed as
\[
    \mathcal{G}^R(\mathbf{k}, \sigma, t; t') = -i \theta(t - t') \langle \{ \hat{c}_{\mathbf{k}, \sigma}(t),\, \hat{c}_{\mathbf{k}, \sigma}^\dagger(t') \} \rangle,
\]
since the Hamiltonian is diagonal in momentum space, and thus the Green's function is also diagonal in momentum space. Let us use this to compute a free electron Green's function.

\begin{example}[Green's function for free electrons]
    Consider a system of free electrons described by the Hamiltonian
    \[
        \hat{H} = \hat{H}_0 = \sum_{\mathbf{k}, \sigma} \xi_{\mathbf{k}} \hat{c}_{\mathbf{k}, \sigma}^\dagger \hat{c}_{\mathbf{k}, \sigma},
    \]
    where \(\xi_{\mathbf{k}}\) is the energy of an electron with momentum \(\mathbf{k}\):
    \[
        \xi_{\mathbf{k}} = \epsilon_{\mathbf{k}, \sigma} - \mu = \frac{\hbar^2 k^2}{2m} - \mu.
    \]
    The Green's function for this system can be computed using the definition:
    \[
        \mathcal{G}^R(\mathbf{k}, \sigma, t; t') = -i \theta(t - t') \langle \{ \hat{c}_{\mathbf{k}, \sigma}(t),\, \hat{c}_{\mathbf{k}, \sigma}^\dagger(t') \} \rangle.
    \]
    Since the Hamiltonian is diagonal in momentum space and we are in the Heisenberg picture, we can express the time evolution of the creation and annihilation operators as
    \[
        \frac{\mathrm{d}}{\mathrm{dt}} \hat{c}_{\mathbf{k}, \sigma}(t) = i \left[ \hat{H}, \hat{c}_{\mathbf{k}, \sigma}(t) \right] = i e^{i \hat{H} t} \left[ \hat{H}, \hat{c}_{\mathbf{k}, \sigma} \right] e^{-i \hat{H} t},
    \]
    where expanding the commutator we get
    \[
        \begin{aligned}
            \left[ \hat{H}, \hat{c}_{\mathbf{k}, \sigma} \right] & = \sum_{\mathbf{k}', \sigma'} \xi_{\mathbf{k}'} \left[ \hat{c}_{\mathbf{k}', \sigma'}^\dagger \hat{c}_{\mathbf{k}', \sigma'},\, \hat{c}_{\mathbf{k}, \sigma} \right]  = -\xi_{\mathbf{k}} \hat{c}_{\mathbf{k}, \sigma},
        \end{aligned}
    \]
    where we have expanded the commutator and used the anticommutation relations of the fermionic operators \([AB,C]=A\{B,C\} - \{A,C\} B\). Thus we have
    \[
        \frac{\mathrm{d}}{\mathrm{dt}} \hat{c}_{\mathbf{k}, \sigma}(t) = -i \xi_{\mathbf{k}} \hat{c}_{\mathbf{k}, \sigma}(t) = -i \xi_{\mathbf{k}} e^{i \hat{H} t} \hat{c}_{\mathbf{k}, \sigma} e^{-i \hat{H} t},
    \]
    showing that the fermionic annihilation operator for a non interacting system evolves in time as the previous expression, and thus we can write the Green's function as\TODO{correct expression.}
    \[
        \begin{aligned}
            \mathcal{G}^R(\mathbf{k}, \sigma, t; t') & = -i \theta(t - t') \langle \{ e^{i \hat{H} t} \hat{c}_{\mathbf{k}, \sigma} e^{-i \hat{H} t},\, e^{i \hat{H} t'} \hat{c}_{\mathbf{k}, \sigma}^\dagger e^{-i \hat{H} t'} \} \rangle \\
                                                     & = -i \theta(t - t') e^{-i \xi_{\mathbf{k}} (t - t')}.
        \end{aligned}
    \]

    We found the Fermi function, since the thermal everage of the number operator \(\hat{n}_{\mathbf{k}, \sigma} = \hat{c}_{\mathbf{k}, \sigma}^\dagger \hat{c}_{\mathbf{k}, \sigma}\) for fermions is given by the Fermi-Dirac distribution:
    \[
        \begin{aligned}
            \langle \hat{c}_{\mathbf{k}, \sigma}^\dagger \hat{c}_{\mathbf{k}, \sigma} \rangle = \frac{1}{Z} \mathrm{Tr} \left( e^{-\beta H} (\hat{c}_{\mathbf{k}, \sigma}^\dagger \hat{c}_{\mathbf{k}, \sigma}) \right) = \frac{1}{Z}  \frac{\sum_{n_{\mathbf{k}}=0,1} \hat{n}_{\mathbf{k}} e^{-\beta \xi_{\mathbf{k}} \hat{n}_{\mathbf{k}}}}{\sum_{n_{\mathbf{k}}=0,1} e^{-\beta \xi_{\mathbf{k}} \hat{n}_{\mathbf{k}}}} = \frac{1}{e^{\beta \xi_{\mathbf{k}}} + 1} = n_F(\xi_{\mathbf{k}}).
        \end{aligned}
    \]

    Going back to the Green's function, we can also express the other thermal average as \(1-\hat{n}_{\mathbf{k}}\), using the anticommutation relations of the fermionic operators, and thus we can write the Green's function of a free electron system as
    \[
        \begin{aligned}
            \mathcal{G}^R(\mathbf{k}, \sigma, t; t') = -i \theta(t - t') e^{-i \xi_{\mathbf{k},\sigma} (t - t')}, \quad \propto e^{-\frac{t}{\tau}}.
        \end{aligned}
    \]
    For an interacting hamiltonian, the Green's function will have a more complicated form, and it will contain information about the interactions and correlations in the system. The longer the time difference \(t - t'\), the more interactions and correlations will be taken into account in the Green's function, and thus it will deviate more from the free electron case.
\end{example}

Thus the idea is to arrive and compute the thermal average for the interacting system, which is a very hard task, and thus we will use perturbation theory to compute the Green's function for the interacting system, starting from the free electron Green's function as the zeroth order approximation. This will allow us to capture the effects of interactions and correlations in a systematic way, and to understand how they modify the properties of the system compared to the non-interacting case.

Instead of time, usually the central quantity of interest is the Fourier transform of the Green's function in frequency space, which is given by
\[
    \mathcal{G}^R(\mathbf{k}, \sigma, \omega) = \int_{-\infty}^{\infty} \mathrm{d} t e^{i \omega t} \mathcal{G}^R(\mathbf{k}, \sigma, t; 0) = \int_{-\infty}^{\infty} \mathrm{d} t e^{i \omega t} (-i) \theta(t) e^{-i \xi_{\mathbf{k}} t} e^{- \delta t},
\]
where we have used the Fourier transform of the Heaviside step function and the exponential function; we set \(t^{\prime} = 0\) since it just corresponds to a shift in the time variable, and thus it does not affect the physics of the problem. Furthermore we added a small imaginary part \(\delta\) to the exponent to ensure convergence of the integral, which is a common technique in many body theory to handle the singularities that can arise in the Green's function. The result of this integral, after the integration of the delta, is
\[
    \mathcal{G}^R(\mathbf{k}, \sigma, \omega) = -i \int_0^{\infty} \mathrm{d} t e^{i (\omega - \xi_{\mathbf{k}, \sigma} + i \delta) t} = -i \left[ \frac{e^{i (\omega - \xi_{\mathbf{k}, \sigma} + i \delta) t}}{i (\omega - \xi_{\mathbf{k}, \sigma} + i \delta)} \right]_0^{\infty} = \frac{1}{\omega - \xi_{\mathbf{k}, \sigma} + i \delta},
\]
where the importance of the small exponential factor \(e^{- \delta t}\) is evident, since it ensures that the integral converges and that the Green's function has the correct analytic properties in the complex frequency plane. The poles of the Green's function in the complex frequency plane correspond to the energy levels of the system, and the imaginary part of the poles gives information about the lifetime of the quasiparticles in the system.

Every time we Fourier transform the Green's function, we are changing the representation of the Green's function from the time domain to the frequency domain, and thus we are changing the way we look at the physics of the problem. The \(\delta\) term in the Green's function is crucial for ensuring that the Green's function has the correct analytic properties: it is a slight shift in the complex frequency plane that ensures that the poles of the Green's function are located in the correct half-plane, which is essential for maintaining causality and for correctly describing the response of the system to external perturbations.

We found \(\mathcal{G}_0^R(\mathbf{k}, \sigma, \omega)\) the retarded Green's function for a free electron system, which is the starting point for our perturbative expansion
\[
    \mathcal{G}_0^R(\mathbf{k}, \sigma, \omega) = \frac{1}{\omega - \xi_{\mathbf{k}, \sigma} + i \delta}.
\]

\subsection{Spectral Function}
Associated to this green's function, we can also define the \textbf{spectral function}, which is given by the imaginary part of the Green's function:
\[
    A (\mathbf{k}, \sigma) = -\frac{1}{\pi} \mathrm{Im} \mathcal{G}_0^R(\mathbf{k}, \sigma, \omega) = - \frac{1}{\pi} \mathrm{Im} \left( \frac{1}{\omega - \xi_{\mathbf{k}, \sigma} + i \delta} \right),
\]
where we know
\[
    \frac{1}{x + i \delta} = \mathcal{P} \left( \frac{1}{x} \right) - i \pi \delta(x),
\]
where \(\mathcal{P}\) denotes the Cauchy principal value, and thus we can write the spectral function as
\[
    A (\mathbf{k}, \sigma) = - \frac{1}{\pi} \mathrm{Im} \left( \frac{1}{\omega - \xi_{\mathbf{k}, \sigma} + i \delta} \right) \frac{\omega - \xi_{\mathbf{k}, \sigma} - i \delta}{\omega - \xi_{\mathbf{k}, \sigma} - i \delta}
\]
where we multiplied and divided by the denominator with a changed sign in front of \(i \delta\), and computing the imaginary part we get
\[
    A(\mathbf{k}, \sigma) = - \frac{1}{\pi} \mathrm{Im} \left( \frac{\omega - \xi_{\mathbf{k}, \sigma} - i \delta}{(\omega - \xi_{\mathbf{k}, \sigma})^2 + \delta^2} \right) = \frac{1}{\pi} \left( \frac{\delta}{(\omega - \xi_{\mathbf{k}, \sigma})^2 + \delta^2} \right) \xrightarrow{\delta \to  0} \frac{1}{\pi} \frac{\delta}{(\omega - \xi_{\mathbf{k}, \sigma})^2 + \delta^2},
\]
since the lorentzian function \(\frac{1}{\pi} \frac{\delta}{x^2 + \delta^2}\) converges to the Dirac delta function \(\delta(x)\) in the limit \(\delta \to 0\), we can write the spectral function as
\[
    A(\mathbf{k}, \sigma) = \delta(\omega - \xi_{\mathbf{k}, \sigma}),
\]
which means that the spectral function for a free electron system is a delta function centered at the energy \(\xi_{\mathbf{k}, \sigma}\). This indicates that the energy levels of the system are sharply defined, and that the quasiparticles have an infinite lifetime (since there is no broadening of the spectral function).

Thus the only possible energy for a particle with momentum \(\mathbf{k}\) and spin \(\sigma\) is \(\xi_{\mathbf{k}, \sigma}\), and the spectral function tells us that the density of states at this energy is infinite, while it is zero at all other energies. This is a characteristic feature of non-interacting systems, where the energy levels are discrete and sharply defined. In contrast, in interacting systems, the spectral function will generally have a more complex structure that reflect the effects of interactions and correlations in the system.

If we now add a term in the Green's function to account for the interactions in the system: a time exponential decay term \(e^{-\frac{t-t^{\prime}}{\Gamma}}\), where \(\Gamma\) takes into account the interactions and correlations in the system, and thus it gives a finite lifetime to the quasiparticles in the system. The Green's function for an interacting system can thus be written as
\[
    \mathcal{G}^R(\mathbf{k}, \sigma, t; t') = -i \theta(t - t') e^{-i \xi_{\mathbf{k}, \sigma} (t - t')} e^{-\frac{t-t^{\prime}}{\Gamma}},
\]
where the exponential decay term \(e^{-\frac{t-t^{\prime}}{\Gamma}}\) captures the effects of interactions and correlations in the system, and it leads to a broadening of the spectral function, which indicates that the energy levels of the system are no longer sharply defined, and that the quasiparticles have a finite lifetime. The Fourier transform of this Green's function will give us a more complex expression for the Green's function in frequency space, which will contain information about the interactions and correlations in the system. The spectral function for this interacting system can be computed as
\[
    A(\mathbf{k}, \sigma) = -\frac{1}{\pi} \mathrm{Im} \mathcal{G}^R(\mathbf{k}, \sigma, \omega) = \frac{1}{\pi} \frac{(1/\Gamma)}{(\omega - \xi_{\mathbf{k}, \sigma})^2 + (1/\Gamma)^2},
\]
where we have used the Fourier transform of the Green's function with the exponential decay term, and we can see how the interactions and correlations in the system lead to a broadening of the spectral function, which indicates that the energy levels of the system are no longer sharply defined, and that the quasiparticles have a finite lifetime. The width of the spectral function is given by \(1/\Gamma\), which is inversely proportional to the lifetime of the quasiparticles: the shorter the lifetime, the broader the spectral function, and vice versa. This illustrates how interactions and correlations in the system can lead to a finite lifetime for quasiparticles, and how this is reflected in the properties of the Green's function and the spectral function.\TODO{Add image of parabula with delta functions and lorentzian functions.}

Thus given the interaction we need a machinery capable of computing \(\Gamma\), in order to obtain the Green's function for the interacting system, and thus to understand the properties of the system in the presence of interactions and correlations. This is where perturbation theory and Feynman diagrams come into play, as they provide a systematic way to compute the effects of interactions on the Green's function and to extract physical quantities such as \(\Gamma\) from the theory. In a certain sense we are measuring how many states are available for the particle to decay into, and thus how many interactions it can have with other particles in the system, which is what gives rise to the finite lifetime of the quasiparticles in the system. Thus we can generalize the Densitiy of States (DeS) to the case of an interacting system, by using the spectral function, which is given by
\[
    \mathrm{DeS}(\omega) = \sum_{\mathbf{k}, \sigma} A(\mathbf{k}, \sigma, \omega),
\]
where the spectral function \(A(\mathbf{k}, \sigma, \omega)\) captures the effects of interactions and correlations in the system, and thus it gives us a more accurate picture of the density of states in the presence of interactions. The density of states for an interacting system will generally have a more complex structure than the non-interacting case, and it will reflect the effects of interactions and correlations in the system, such as the broadening of energy levels and the emergence of new features in the density of states due to interactions.

Experiments in angular resolved photoemission spectroscopy (ARPES) can directly measure the spectral function \(A(\mathbf{k}, \sigma, \omega)\) of a material, providing insights into the electronic structure and interactions in the system. By analyzing the ARPES data, we can extract information about the energy dispersion, the lifetime of quasiparticles, and the effects of interactions and correlations in the material. This makes the spectral function a crucial quantity for understanding the physics of condensed matter systems and for interpreting experimental results.
